\documentclass{article}
\usepackage{amsmath,amssymb}
\usepackage{xurl}
\usepackage[hidelinks]{hyperref}
\setlength{\emergencystretch}{2em}
\widowpenalty=10000
\clubpenalty=10000
% Shared commands for notes in docs/paper-gaps/.

\DeclareUnicodeCharacter{2080}{\ifmmode{}_{0}\else\textsubscript{0}\fi}
\DeclareUnicodeCharacter{2081}{\ifmmode{}_{1}\else\textsubscript{1}\fi}
\DeclareUnicodeCharacter{2082}{\ifmmode{}_{2}\else\textsubscript{2}\fi}
\DeclareUnicodeCharacter{2099}{\ifmmode{}_{n}\else\textsubscript{n}\fi}

\newcommand{\Lean}{\textsc{Lean}}
\ExplSyntaxOn
\NewDocumentCommand{\leanid}{m}{
  \ifmmode
    \textsf{\detokenize{#1}}
  \else
    \group_begin:
    \urlstyle{sf}
    \tl_set:Nn \l_tmpa_tl {#1}
    \tl_replace_all:Nnn \l_tmpa_tl {\_} {_}
    \exp_args:NV \path \l_tmpa_tl
    \group_end:
  \fi
}
\ExplSyntaxOff
\providecommand{\tr}{\operatorname{tr}}
\newcommand{\tnleanIssueBase}{https://github.com/LionSR/TNLean/issues/}
\newcommand{\tnleanPrBase}{https://github.com/LionSR/TNLean/pull/}
\newcommand{\tnleanDocBase}{https://sirui-lu.com/TNLean/docs/}
\newcommand{\ghissue}[1]{\href{\tnleanIssueBase#1}{GitHub issue \##1}}
\newcommand{\ghpr}[1]{\href{\tnleanPrBase#1}{PR \##1}}
\newcommand{\doclink}[1]{\href{\tnleanDocBase#1.html}{\path{#1}}}

% Dirac notation and the identity operator, as in blueprint/src/macros/common.tex.
% \providecommand keeps note-local definitions authoritative.
\usepackage{dsfont}
\providecommand{\ket}[1]{|#1\rangle}
\providecommand{\bra}[1]{\langle #1|}
\providecommand{\braket}[2]{\langle #1 | #2 \rangle}
\providecommand{\ketbra}[2]{|#1\rangle\!\langle #2|}
\providecommand{\Id}{\mathds{1}}
\providecommand{\one}{\mathds{1}}
\providecommand{\C}{\mathbb{C}}
\providecommand{\MN}[1]{M_{#1}(\C)}
\providecommand{\MD}{M_D(\C)}

% Verdict marker: \gapnote{<kind>}{<status>}. Kind is the policy
% classification (clarification, local-correction, scope-restriction,
% unfaithful, false-source, open-gap); status is open, wip (elimination
% underway), resolved, or historical. Machine-read by the site index;
% prints a verdict line.
\newcommand{\gapnote}[2]{\par\noindent\textbf{Verdict:} #1 (#2).\par}

\title{SPT Fixed Points with a Nontrivial Character}
\author{}
\date{2026-10-03}
\begin{document}
\maketitle
\gapnote{scope-restriction}{resolved}

\section{The source statement}
In Section~III.A of \cite{Cirac2021Matrix}, a symmetry acts on an MPS by
$S_g(A^i)=e^{i\theta(g)}X_g^\dagger A^iX_g$, where
$\phi(g)=e^{i\theta(g)}$ is a circle-valued character. The source constructs
a zero-correlation-length representative for a prescribed character and
2-cocycle class. On a ring of $N$ sites, the character contributes the
physical phase $\phi(g)^N$. The relevant passage is
\path{Papers/2011.12127/TN-Review-main.tex}, lines 1147--1157.

\section{The earlier restriction}
The original exact-invariance predicate required each physical twist to
have the same matrix product vectors as the untwisted tensor. Consequently,
the fixed-point symmetry and class-realization results using that predicate
were restricted to $\phi=1$, although the local twist identity was already
proved for arbitrary $\phi$. Exact vector invariance omits the character
phase when $\phi(g)^N\ne1$.

\section{Character and cocycle-class realization}
Let $\rho$ be a projective representation with factor system $\omega$,
put $W_g=\rho(g^{-1})$, and take the dimer fixed point
\begin{align}
    A^{(a,b)} & =D^{-1/2}|a\rangle\langle b|, &
    U(g) & =\phi(g)(W_g^T\otimes W_g^{-1}).
    \label{eq:rmp_spt_character_action}
\end{align}
The symmetry relation is
\begin{align}
    \widetilde A_g^i=\sum_jU(g)_{ij}A^j
      & =\phi(g)W_gA^iW_g^{-1}.
    \label{eq:rmp_spt_trivial_character_twist}
\end{align}
Thus the twist is gauge equivalent to $\phi(g)A$, and
\begin{align}
    |V^{(N)}(\widetilde A_g)\rangle
      & =\phi(g)^N|V^{(N)}(A)\rangle.
    \label{eq:rmp_spt_character_vector}
\end{align}
These statements use on-site symmetry up to a character, rather than
requiring equality of the two vectors. The projective representation
$\rho$ itself realizes (\ref{eq:rmp_spt_trivial_character_twist}).
For $D\geq1$, every other virtual projective representation realizing
that relation with the same character has a factor system cohomologous
to $\omega$: cancel the nonzero scalar $\phi(g)$, then apply gauge
uniqueness to the injective fixed-point tensor.

For finite $G$, a unitary twisted regular representation supplies the
virtual representation from the cocycle alone. Its factor system is
cohomologous to $\omega$, and the physical action in
(\ref{eq:rmp_spt_character_action}) realizes the prescribed character and
cocycle class. No unit-modulus assumption on a complex-valued character
is needed here. Every $g$ has finite positive order $m$, so
$\phi(g)^m=1$ and $|\phi(g)|=1$ follow from the character law. For a general
group, the algebraic identities hold for every complex-valued character;
unitarity of the physical action requires the source's circle-valued
character and unitary virtual matrices.

The general-character results are in
\doclink{TNLean/MPS/Symmetry/Character} and
\doclink{TNLean/MPS/Symmetry/SPTFixedPoint}. They include the vector phase,
on-site symmetry up to the character, virtual class uniqueness and
existence, and finite-group realization from the cocycle alone. The
trivial-character restriction is therefore resolved.

\section{Distinct source questions}
This correction concerns the character phase. The use of the dimer fixed
point in place of the printed tensor remains the local correction recorded
in \cite{gap:rmp_spt_fixed_point_tensor}. The present action treats linear
on-site symmetries; time reversal and reflection, with their twisted
cocycles, remain outside this result as recorded in
\cite{gap:rmp_spt_fixed_point_onsite_scope}. Resolving the character
restriction does not alter either of these separate conclusions.

\bibliographystyle{alpha}
\bibliography{references}
\end{document}
